\subsection{Paths Forward} \label{sec:paths}

The original position paper, written over a year ago, discusses a collection of paths forward, including richer code data, training in executable 
environments, and adaptation to specialized codebases. In this thesis, I will focus on inference-time 
approaches: improving retrieval over code, combining neural models with symbolic program analysis, 
and integrating software-engineering tools into model workflows.
Because the field of AI for code has moved rapidly over the past year, there has already 
been significant progress on several of these challenges, validating the importance of several of
these directions.

\subsubsection{Inference-Time Approaches}

\par\medskip\noindent\textbf{Semantic-Aware Embeddings and Retrieval}\phantomsection\label{sec:directions-retrieval}\par\smallskip

\textbf{Semantic and execution-aware code embeddings}: A year ago, code embeddings often treat programs as 
pure token sequences rather than explicitly incorporating execution and semantics. 
While not discussed in this thesis, something we discovered throughout careful analysis was that
code that is close in embedding space is often syntactically similar rather than semantically similar \citep{utpala2023language}. 
Before the LLM era, however, there were efforts to learn program representations that encode semantics, 
including neural modules for program operations \citep{nye2020representing} and 
execution-aware latent representations for partial or full programs \citep{zohar2018automatic, ellis2019write, chen2021latent}. 
While today's language models reason about semantics beyond just syntax, they still use syntactic similarity as a strong
signal for retrieval.  
Revisiting these ideas may improve both code retrieval abilities while making code 
representations more aligned with program behavior.

\textbf{Better retrieval-augmented code generation}: Retrieval is especially useful in 
specialized contexts where the model has not seen enough examples to fully understand a 
language, API, or codebase. Documentation, function definitions, and local usage examples 
can teach syntax, available abstractions, and conventions. However, retrieving relevant 
code is only half the problem. The model must also reuse the retrieval correctly, 
adapting algorithms, translating between languages, and recognizing subtle differences 
between the retrieved context and the current task. This suggests that retrieval-augmented 
coding agents should be designed not only to consume retrieved snippets, but also to critically 
reason about how those snippets should be modified. One failure mode of today's systems is that they 
often copy retrieved code patterns without modification, even when the new context that the
model is working in has subtle differences that require adaptation. 


\textbf{Retrieving via code navigation on the fly}: Static retrieval indexes can be expensive to build and can become stale as a codebase evolves. 
A year ago, we suggested that instead of relying only on static embeddings, agents should retrieve information by 
navigating the codebase directly, using commands such as \texttt{cd}, \texttt{ls}, and \texttt{grep}, 
IDE actions such as jumping to definitions or finding references, and static-analysis tools that expose ASTs or 
dependency graphs. The community has made significant progress in this direction. In fact, it was
a significant unlock that made coding agents drastically more usable. Searching and navigating files
on the fly, in addition to exploring the project filesystem in other ways, is now a core aspect of Claude Code and Codex.

At the same time, there is still significant headroom in how agents retrieve information from a codebase. 
Future systems could use more sophisticated search procedures that explicitly exploit program structure,
such as dependency graphs, call graphs, or ASTs, rather than treating the filesystem as a flat collection of files. 
They could also make the search process faster, since current agents often spend a large fraction of their time 
exploring before they can make progress on the actual task. Finally, as context windows become longer, agents may be able 
to learn from their own retrieval histories: instead of treating each search as independent, they could track which files, 
commands, and retrieved snippets were actually useful and use that information to guide future navigation.

\par\medskip\noindent\textbf{Neurosymbolic Approaches}\phantomsection\label{d:sec-neurosymbolic}\par\smallskip

Code is a unique domain because there is a vast body of programming-languages research to build on, including abstract interpretation \citep{cousot1977abstract}, concolic testing \citep{godefroid2005dart, sen2005cute}, model checking \citep{clarke1997model}, linting, and type-checking \citep{cardelli1996type}. These tools can surface bugs, enforce invariants, narrow search spaces, and provide structured feedback. Traditional PL tools often require precise specifications, can be computationally expensive, and may produce false positives. LLMs can help mitigate these issues by proposing specifications, localizing likely errors, and deciding when a symbolic tool is useful.

LINC and ProofOptimizer highlight this broader pattern. In LINC, the model
translated natural language reasoning into executable symbolic programs. In ProofOptimizer,
the model proposed proof transformations while the compiler ensured correctness.
Neurosymbolic methods potentially allow AI systems to explore a large search space while maintaining guarantees throughout.

There are several ways this could potentially be integrated. During code generation, 
program-analysis techniques could be applied to shorter snippets of AI-generated code to 
surface bugs or prove properties. To improve code understanding, agents can inspect program 
structure such as abstract syntax trees \citep{gong2024ast} in the same way that they navigate
the codebase and filesystem today. 
When debugging a large codebase, AI could first narrow down the suspicious region before handing 
the problem to more precise symbolic tools. During code generation in DSLs, LLMs can leverage grammars 
for constrained decoding~\citep{poesia2022synchromesh, geng2023grammar, wei2023copiloting}. During refactoring, 
static analysis can identify whether new errors have been introduced and cut off unpromising 
search paths.

\par\medskip\noindent\textbf{Incorporating SWE Tools}\phantomsection\label{d:sec-agents-tools}\par\smallskip

Software engineers integrate a variety of domain-specific tools when writing code. AI systems should similarly be able to invoke tools during inference and incorporate their outputs into the next step of reasoning. At inference time, this requires three skills: choosing the right tool, invoking it correctly, and incorporating its output into the next step of reasoning. This requires precise tool interfaces, reliable access to execution or analysis results, and agent workflows that can use tool feedback without losing track of the broader task. Useful tools include compilers, test runners, linters, type checkers, package managers, profilers, debuggers, and build systems; the core challenge is not simply exposing these tools to a model, but teaching the model when their feedback is relevant and how to act on it.

Taken together, these inference-time directions suggest a common principle: the next stage of AI for code should be organized around feedback-rich interaction with programs. Retrieval helps models understand local codebases, neurosymbolic approaches connect model reasoning to program semantics, and tool use grounds model reasoning in executable evidence.
